\chapter{2024 Nobel Prize in Economics}
\textbf{Laureates: }Daron Acemoglu (Turkey/USA), Simon Johnson (UK), James Robinson (USA)

\section*{Reason for Award}
For their research on how institutions are formed and affect prosperity

\section{What They Did}
Daron Acemoglu, Simon Johnson, and James Robinson developed comparative
institutional analysis, showing how institutions determine national wealth.

They distinguished between inclusive institutions, which secure property rights,
enforce contracts, and provide public services, and extractive institutions,
which concentrate power and wealth in the hands of a few. Inclusive institutions
create incentives for investment, innovation, and participation. Extractive
institutions create incentives for rent-seeking and appropriation.

Their research on the colonial origins of comparative development showed that
colonization strategies created different institutional trajectories. Settler
colonies with favorable disease environments developed inclusive institutions
because settlers demanded property rights and representative government.
Extractive colonies developed institutions designed to extract resources for
the metropole. These institutional differences persisted after independence,
explaining current income disparities.

Johnson, Acemoglu, and Robinson showed that geography, culture, and religion
are not the fundamental causes of development differences. Instead, political
and economic institutions shape incentives and determine outcomes. Their
instrumental variables approach used European settlement mortality rates as
an instrument for current institutional quality, providing credible causal
evidence.

They documented how geographic, cultural, and political factors interact with
institutions to shape economic outcomes. Their work emphasized that institutions
are not destiny but can change through political struggle and policy reform.

Their research on the Reversion to Meaninglessness argued that institutional
improvements in one country do not automatically spread to others. Each
country must develop its own inclusive institutions through internal political
processes.

\section{Impact on Economic Thinking}
Acemoglu, Johnson, and Robinson transformed development economics and political
economy. Their institutional approach challenged conventional explanations based
on geography, culture, and ignorance.

They showed that differences in institutions explain most variation in income
per capita across countries. The quantitative impact of institutions on income
is large, with institutional differences accounting for most of the income gap
between rich and poor nations.

Their research influenced policy debates about governance, rule of law, and
institutional reform. The concept of extractive versus inclusive institutions
provided a framework for understanding why some countries develop while others
stagnate.

Their work inspired research on political economy, transition economies, and
institutional change. Scholars examined how institutions evolve, how elites
maintain extractive systems, and how inclusive institutions can be established.

The emphasis on institutions as the fundamental cause of prosperity reshaped
development economics and policy thinking. International organizations
incorporated institutional quality metrics into their assessments and lending
criteria.

\section{Modern Perspectives Today}
Institutional analysis based on Acemoglu, Johnson, and Robinson's work remains
central to development economics and political science.

Research on governance, corruption, and state capacity builds on their framework.
Studies examine how institutional quality affects economic performance, social
outcomes, and political stability.

Studies of democratic transitions, authoritarian resilience, and institutional
reform draw on their insights. Research has explored how elites maintain power,
how social movements can drive institutional change, and how constitutional
design affects outcomes.

The role of institutions in economic growth, innovation, and inequality continues
to be explored. Recent work examines how digital technologies, climate change,
and globalization interact with institutions.

The concept of inclusive versus extractive institutions informs policy debates
about development assistance, conditionality, and institutional reform. As
countries grapple with inequality, polarization, and democratic backsliding,
their institutional framework remains essential for understanding these
challenges.

\section*{Important Bibliography}
\begin{itemize}
\item Acemoglu, D., Johnson, S., \& Robinson, J. (2001). ``The colonial origins of comparative development: An empirical investigation.'' \textit{American Economic Review}, 91(5), 1369--1401.
\item Acemoglu, D., \& Robinson, J. (2012). \textit{Why Nations Fail: The Origins of Power, Prosperity, and Poverty}. Crown Publishers.
\item Acemoglu, D., Johnson, S., \& Robinson, J. (2005). `` Institutions as the fundamental cause of long-run growth.'' \textit{Handbook of Economic Growth}, 1, 385--472.
\item Acemoglu, D. (2024). ``Daron Acemoglu: The 2024 Nobel Memorial Prize in Economic Sciences.'' \textit{Journal of Economic Perspectives}, 38(4), 1--12.
\item Johnson, S. (2024). ``Simon Johnson: The 2024 Nobel Memorial Prize in Economic Sciences.'' \textit{Journal of Economic Perspectives}, 38(4), 13--24.
\end{itemize}
